

This work demonstrates the feasibility of using steering vectors to guide model responses toward specific answers or answer types. Ideally, steering should be precisely targeted, so that a shift from a “yes” to a “no” affects only questions initially answered with “yes”. While this effect is often achieved, it is not perfect and consistent across all cases. Our experiments focus on the VQA, which involves relatively short answers. We also explore an extension to image captioning tasks, showing that captions can be modified to include specific stylistic elements, such as sentiment or color-related words. Targeted steering in tasks with longer answers proves more challenging. We also notice a tradeoff between the steering strength and the quality of generated responses. For instance, with longer outputs, steering may increase the likelihood of hallucinated content. This approach can be extended to other applications, such as directing outputs toward specific images or questions, addressing biases, or mitigating safety concerns.

testing different ways to compute steering vectors



